\documentclass[conference]{IEEEtran}

% --- Grundkonfiguration ---
\usepackage[utf8]{inputenc}
\usepackage[T1]{fontenc}
\usepackage[ngerman]{babel}

% --- Wissenschaftliche Pakete ---
\usepackage{cite}
\usepackage{amsmath,amssymb,amsfonts}
\usepackage[demo]{graphicx} % [demo] verhindert pdflatex-Absturz bei fehlerhaften Bilddaten
\usepackage{url}
\usepackage{booktabs} 
\usepackage{pgfplotstable} 
\pgfplotsset{compat=1.18}

% --- Pfaddefinitionen (Anpassung an project/expose/) ---
% Geht zwei Ebenen hoch zu ~/ITS_THB/Hardware-Sicherheit/assets/
\graphicspath{{../../assets/}}

\begin{document}

\title{Exposé: Analyse von Hardware-Virtualisierung unter Proxmox VE}

\author{\IEEEauthorblockN{Maximilian Wagner}
\IEEEauthorblockA{\textit{Studiengang IT-Sicherheit} \\
\textit{Technische Hochschule Brandenburg}\\
maximilian.wagner@th-brandenburg.de}
}

\maketitle

\begin{abstract}
Dieses Dokument skizziert die methodische Untersuchung von Virtualisierungs-Overheads und Hardware-Isolationsmechanismen.
\end{abstract}

\section{Einleitung}
Im Rahmen dieses Projekts wird die Effizienz der Intel VT-x Erweiterungen evaluiert \cite{intelvtx}. Der Fokus liegt auf der Quantifizierung von VM-Exit-Latenzen.

\section{Experimenteller Aufbau}
Das Test-Setup basiert auf einer Debian-Gastinstanz innerhalb eines Proxmox-Hypervisors. Abbildung \ref{fig:setup} verdeutlicht die Architektur.

\begin{figure}[htbp]
    \centering
    \includegraphics[width=0.45\textwidth]{picture.jpg}
    \caption{Architekturdiagramm der Testumgebung.}
    \label{fig:setup}
\end{figure}

\section{Benchmark-Datenstruktur}
Die Ergebnisse werden automatisiert in CSV-Dateien überführt. Tabelle \ref{tab:results} zeigt das erwartete Datenformat.

\begin{table}[htbp]
\caption{Beispielhafte Messdatenstruktur}
\centering
% Pfad angepasst: Eine Ebene hoch, dann in experiments/
% Spalten-Aliase definieren, um Math-Mode-Fehler durch Unterstriche zu eliminieren
\pgfplotstabletypeset[
    col sep=semicolon, 
    string type,
    columns/CPU_Events_per_sec/.style={column name={CPU (Events/s)}},
    columns/IOPS_Random_Write/.style={column name={I/O (IOPS)}},
    every head row/.style={before row=\toprule,after row=\midrule},
    every last row/.style={after row=\bottomrule},
]{../experiments/summary.csv} 
\label{tab:results}
\end{table}

\bibliographystyle{IEEEtran}
\bibliography{references_expose} % Name der lokalen Bib-Datei

\end{document}